\chapter[Introduction]{Introduction}
\label{chap:introduction}

\lettrine[
  lines=4,
]{\color{OxfordBlue} S}
{ince} antiquity evolutionary strategies have been employed to solve engineering problems. In the 26\textsuperscript{th} century BC, the builders of the ancient Egyptian Bent pyramid perturbed the slope of the sides of the pyramid from $55^\circ$ to $43^\circ$ halfway through construction. This allowed each new stone course to sit further inward, placing less weight on the lower structure. In the subsequent Red Pyramid, they used the same safer $43^\circ$ slope from the foundations upward~\citep{lehner1997pyramids}.

Although the builders did not describe their work mathematically, their approach followed a ubiquitous logic: propose a deisgn, test its performance, and revise based only on the observed outcome. This is precisely the setting of zeroth-order optimization, with acess to an oracle $f(x)$ but not its gradient $\grad f(x)$. Gradient-free optimization is the natural setting for problems where derivatives are unavailable or prohibitively expensive, including energy systems, materials science, computational engineering design~\citep{alarie2021blackbox}, automated experimental design~\citep{shahriari2016human}, hyperparameter selection~\citep{bergstra2012random}, robotics and reinforcement learning~\citep{cully2015robots,salimans2017evolution}, molecular design~\citep{gomez2018automatic}, and black-box attacks on neural networks~\citep{chen2017zoo}.

Optimization algorithms which access $f(x)$ but not $\grad f(x)$ are called zeroth-order. The first strategy one might try is to evaluate at a nearby point $x+\sigma p$ for a perturbation $p \in \R^d$, accept it if it is better, and reject it otherwise. For Gaussian perturbations, however, the best alignment among $r$ candidates grows only as $\sqrt{\log r}$, yielding diminishing returns in the population size. Evolutionary strategies instead use all $r$ function evaluations to construct the gradient estimator $\frac{1}{\sigma r} \sum_{i=1}^r f(x+\sigma p_i) p_i$, where $\sigma > 0$ is the smoothing radius. This estimator has expectation $\grad f_\sigma(x)$, the gradient of the Gaussian-smoothed objective $f_\sigma(x) = \mbb E_{p \sim \mathcal N(0,I_d)}[f(x+\sigma p)]$, while averaging independent perturbations reduces its variance by a factor of $r$. When function evaluations are performed in parallel, this variance reduction can yield a linear improvement in wall-clock complexity with population size $r$.

Typically, zeroth-order algorithms require $d$ times more function evaluations than their first-order counterparts~\citep{nesterov2017random}. This dimension overhead becomes prohibitive for modern LLMs with billions or trillions of parameters. A Gaussian search perturbation $p \in \R^d$ is dense almost surely, so explicitly constructing and applying it requires $\Theta(d)$ arithmetic. This fundamentally limits how many parallel workers are possible, previously being around 1200~\citep{salimans2017evolution}.

This problem was recently addressed by \citet{sarkar2025evolution} with Evolution Guided GeneRal Optimization via Low-rank Learning (EGGROLL). By taking low-rank perturbations $P=ab^\top$ for $a \sim \mathcal N(0,I_d)$ and $b \sim \mathcal N(0, I_m)$, \citet{sarkar2025evolution} were able to scale zeroth-order optimization to millions of parallel workers. Instead of $\Theta(dm)$ cost for each perturbation, it was vastly reduced to $\Theta(d+m)$ by exploiting $(W + \sigma ab^\top)x = Wx + \sigma a(b^\top x)$. For the first time ever, evolutionary strategies were deployed at the hyperscale.

However it remained open whether this low-rank engineering trick enjoys the same theoretical properties as Gaussian perturbations. We say $x$ is an $\ve$-first-order stationary point of $f \in \mathcal C^1(\R^d)$ if $\mg{\grad f(x)} \leq \ve$. \citet{ren2023escaping} analyze a closely related zeroth-order algorithm using dense Gaussian perturbations. Under some smoothness assumptions, they show that the algorithm converges to an $\ve$-second-order stationary point with high probability in $\widetilde{O}(d/r \cdot 1/\ve^{5/2})$ iterations, with each iteration using $2r$ function evaluations. For the weaker objective of first-order stationarity, the corresponding complexity is $\widetilde{O}(d/r \cdot 1/\ve^2)$. This motivates our first research question:

\begin{quote}
    \textbf{Research Question 1.}
    Does running EGGROLL with its low-rank perturbations converge, and if so, does it match the iteration complexity obtained using dense Gaussian perturbations? 
\end{quote}

The second part of this dissertation considers whether zeroth-order methods can also inherit acceleration. For smooth convex optimization, gradient descent requires $O(1/\ve)$ iterations to reach an $\ve$-optimal point, while Nesterov's Accelerated Gradient Descent improves this to $O(1/\sqrt{\ve})$~\citep{nesterov1983method}. \citet{allenZhu2016linear} interpret this acceleration as a linear coupling of gradient and mirror descent. In the nonconvex setting, \citet{carmon2018accelerated} use accelerated optimization as a sub-routine for sub-quadratic $\widetilde{O}(1/\ve^{7/4})$ convergence to a second-order stationary point, improving on $O(1/\ve^2)$ for gradient desecent. This suggests zeroth-order acceleration could eventually enable improved rates for zeroth-order nonconvex optimization. As a first step, we ask:
\begin{quote}
\textbf{Research Question 2.}
Can zeroth-order linear coupling be analyzed with high probability while retaining its accelerated dependence on $\ve$?
\end{quote}

\section{Contributions}
We make the following contributions.
\begin{itemize}
    \item In Chapter~\ref{chap:fosp-eggroll} we show that, for an $L$-smooth possibly nonconvex function $f$, with probability at least $1-\delta$, (symmetric) EGGROLL with $2r$ function evaluations per iteration converges to a first-order stationary point in 
    \begin{align*}
        \widetilde{O}\left(\frac{dm}{r} \cdot \frac{L(f(X_0) - f^*)}{\ve^2}\right)
    \end{align*}
    iterations, matching dense Gaussian perturbations, and also maintaining the property that parallelizing function evaluations across $r$ workers yields a linear-in-$r$ speedup. In particular, Theorem~\ref{thm:error_concentration} and its consequences yield the general tools needed to generalize many zeroth-order algorithms to the low-rank perturbations case.

    \item In Chapter~\ref{chap:zo-mirror-descent} we show that, with probability at least $1-\delta$, zeroth-order mirror descent with dense Gaussian perturbations, a subroutine for the linear coupling version of accelerated gradient descent, finds an $\ve$-optimal point in 
    \begin{align*}
        \widetilde{O}\left(\frac{d \Theta \rho^2}{\ve^2} \right)
    \end{align*} iterations for convex $\rho$-Lipschitz quadratic functions, matching the rate presented in~\citet{allenZhu2016linear} up to the zeroth-order dimension tax and irrelevant polylogarithmic factors. The analytical techniques developed in this chapter will be critical for analyzing zeroth-order accelerated gradient descent with high probability.

    \item In Chapter~\ref{chap:linear-mixing}, we design a zeroth-order accelerated gradient descent algorithm with dense Gaussian perturbations, showing that for an $L$-smooth convex function $f$, with probability at least $1-\delta$, as long as $r = \widetilde{\Theta}(1)$, it finds an $\ve$-optimal point in
    \begin{align*}
        \widetilde{O}\left(\frac{d}{r} \cdot \sqrt{\frac{L\Theta}{\ve}}\right)
    \end{align*}
    iterations, where we use $2r$ function evaluations per iteration, matching the optimal accelerated rate, with the usual zeroth-order dimension penalty. 
\end{itemize}
